% Chapter 6

\chapter{CONCLUSION AND FUTURE WORK}\label{cha:conclusion}

\section{CONCLUSION}

This project set out to let the payload decide the route. Existing multi-objective UAV
planners build good Pareto fronts and then choose from them with a rule that knows
nothing about the cargo, or with weights an expert re-tunes for every mission. We
separated the two problems. Tier~1 builds an exact NAMOA* front over a
three-dimensional, wind-aware airspace lattice once per region and weather window, and
stores it until it goes stale. Tier~2 turns four payload
tags into objective weights with an explicit formula, scores every route on the front
and returns the best one, or a HOLD when the weather or the battery rules the mission
out. The chosen route has been flown in SITL and in a Gazebo city world.

The preliminary results have been quite encouraging, but there are still some important
observations to make. First, the Tier~1 front was exact for each of the benchmarks, while MOPSO managed to discover 38 routes
out of 115 on the biggest one. On that front the four reference payloads received four
different routes, three of them non-extreme, in about 0.14\,ms each. Secondly, when only the wind direction was varied at constant speed,
the energy consumed on the route that had the lowest cost changed by 27.5\%. The results
of the SITL simulations showed that the drone stayed within 0.13\,m RMS of the planned
trajectory. However, there were a number of flaws found as well. In a steady wind the front
shrank and the presets fell onto two routes; MOPSO ran at reduced settings; the ablation
used fixture data; and the risk weights and battery constants are still placeholders.

\section{FUTURE WORK}

The following experiment conducts the full benchmark (four graphs, nine wind cases,
MOPSO set to use 100 particles, 200 iterations, and 30 runs) with stats for each run and
significance test. Density tags and monotonicity classification need fixing, and our
goal is to find out the reason for Pareto collapse in constant wind case by testing
other normalization and weight multipliers. Ablation and Tier~2 experiments will be
performed again but using real map data, and then Tier~2 will be tested against
crowding distance, knee point, shortest route, and fixed weights for all 225 tag
combinations. There are two flaws in the dashboard that are not fixed yet, and Tier~1
needs better performance in large graphs. Simulation involves putting the Gazebo world
through wind, tuning the battery model according to the class of load it carries, and
running the pre-defined setup for all different wind conditions with full flight logs
recorded. For Phase~II, the system switches to using an F450 quadcopter with the Pixhawk and a
companion computer called Raspberry Pi~5 and a servo for releasing the payload. The
Tier~2 system runs on the drone, and outdoor testing will involve checking the extent to
which it stays on course and the time between entering tags and starting the mission.
Following this, the two-tier system can be further expanded to include multiple drones
having a common front, performing multi-drop delivery missions, or in-flight
re-planning of routes based on the changing weather.
